\chapter{Claim audit: what is settled, and what was withdrawn}
\label{ch:status}
\label{ch:corrections}

The experimental record and the correction record, in one place. They were two
chapters, which meant a claim's current status lived in one and the reason it
changed lived in the other, and checking a number meant knowing to look twice.

Everything below comes from a single configuration, defined once in
\codefile{experiments/frozen\_config.py} and imported by every experiment. The
previous version of this record stitched together runs at three EM iteration
budgets, three replicate counts and two values of $\corr$; several conclusions
turned out to be artefacts of those settings rather than of the model.

\section*{The ledger}
\addcontentsline{toc}{section}{The ledger}

One line per claim. \textbf{Settled} means it can be cited without qualification;
\textbf{conditional} means it holds under a stated restriction that must travel
with it; \textbf{censored} means the run hit a budget rather than a conclusion;
\textbf{withdrawn} means it was claimed and is not any more.

\begin{center}\small
\begin{tabular}{@{}p{0.30\textwidth}p{0.13\textwidth}p{0.47\textwidth}@{}}
\toprule
claim & status & where it stands \\
\midrule
Grid BP reproduces the Gaussian closed form & settled &
  $9.2\times10^{-15}$ at $\ngrid=401$, $\halfwidth=8$; a configuration-specific
  agreement, not a bound for arbitrary kernels \\
Truncation and quadrature are separately controlled & settled &
  \S\ref{sec:rec-rung0}; the boundary statistic is a quantile over chains, not a bound \\
Moment closure $=$ LMMSE of the covariance-matched Gaussian & settled &
  proved, not observed; the gap measures non-Gaussianity \\
Fisher's identity gives the exact gradient & settled &
  $\sim10^{-9}$ against finite differences of the discretised evidence \\
EM ascends the exact evidence & settled &
  monotonicity violation exactly $0$ across every reported run \\
$\corr$ and $\ivar$ are recovered from noised sequences & settled &
  \S\ref{sec:rec-rung1} \\
The innovation \emph{density} is recovered & conditional &
  under the declared density metric and at a certified budget; the shape needs
  $\sim$120 iterations where $\corr$ needs $\sim$25 \\
Structure is worth an order of magnitude against networks & conditional &
  \S\ref{sec:rec-efficiency}; both budgets validation-selected, and the arms are
  asymmetric in supplied information by construction \\
Marginals are blind to the ring's rotation & settled &
  theorem, with a machine-precision zero as its own control \\
Non-Markov violations are not interchangeable & conditional &
  two violations probed, not a general theory \\
Mixture capacity saturates near eight components & withdrawn &
  \S\ref{sec:corr-capacity}; most shape coordinates had not settled at the
  budget used \\
Pointwise and generative fidelity dissociate & withdrawn &
  rests on the same unconverged fits \\
The channel destroys shape recovery & withdrawn &
  a rate at a fixed iteration count, mistaken for an information penalty \\
\bottomrule
\end{tabular}
\end{center}


This chapter is the experimental record. It replaces an earlier version that reported results
from runs at three different EM iteration budgets, three different replicate counts and two
values of $\corr$ --- numbers that could not go in one table, and several of which turned out to
be artefacts of those settings rather than of the model. Part~\ref{part:audit-corrections} of this chapter records what those were and how they
were caught.

Everything below comes from a single configuration, defined once in
\codefile{experiments/frozen\_config.py} and imported by every experiment.

\section{The configuration}

\begin{center}\small
\begin{tabular}{@{}ll@{}}
\toprule
$\corr = 0.85$ & one value everywhere; the earlier $0.8$ runs are retired \\
$\nsites = 32$ & chain length, which is also the ambient dimension \\
$\ngrid = 401$, $\halfwidth = 8$ & read off the grid/domain sweep, \S\ref{sec:rec-rung0} \\
twelve noise levels & log-spaced on $[0.05, 3.0]$ \\
sixteen replicates & everywhere \\
$\nseq \in \{32,\dots,4096\}$ & seven doublings \\
$256$ held-out sequences & drawn from a seed disjoint from every training seed \\
EM stopping & relative $10^{-9}$ on the log-evidence, cap $400$; \textbf{never} on $\corr$ \\
ring & $T = 12$, so $\nsites = 24$ \\
\bottomrule
\end{tabular}
\end{center}

\begin{pitfall}
A frozen configuration is only as good as its reach, and this one has now failed that test twice.
Four experiments carried private replicate knobs that never read \texttt{n\_seeds}, and every one
of them reported \texttt{is\_frozen: true} while running at three, four or eight replicates
(\S\ref{sec:corr-replicates}). The same disease then turned up in the iteration budget:
\texttt{em\_max\_iters} and \texttt{em\_shape\_tol} are declared in the config and read by
\emph{nothing}, while \codefile{exp\_07} passed a literal $120$
(\S\ref{sec:corr-capacity}). Both were invisible to the provenance stamp, which compares the
config against itself. \codefile{tools/check\_replicates.py} now parses the settings dictionaries
for both classes of knob.

The general lesson is worth stating once: a frozen configuration is a claim about what the
experiments read, and only a reader-side check can test that claim. Declaring a value is not the
same as consuming it.
\end{pitfall}

\section{Settled: the numerical foundation}
\label{sec:rec-rung0}

\paragraph{The discretisation is controlled on two axes, separately.} Truncation --- mass outside
$[-\halfwidth,\halfwidth]$ discarded rather than transported --- responds to the domain and not
the resolution: the worst-chain boundary mass falls from $2.3\times10^{-3}$ at $\halfwidth=4$ to
$1.2\times10^{-6}$ at $\halfwidth=8$, essentially independently of $\ngrid$ (the $90$th percentile
falls $2.7\times10^{-4}\to1.5\times10^{-8}$ over the same range; \S\ref{sec:truncation} explains
why both are quoted). Quadrature error falls by a factor of
four per halving of $\gridh$, the predicted second order for the trapezoidal rule. Neither sweep
alone certifies the other, which is why both are reported.

\paragraph{The pipeline is verified where the answer is known.} At $\ngrid=401$, $\halfwidth=8$
the discretised recursion reproduces the Gaussian closed form to $9.2\times10^{-15}$ relative,
and agrees with brute-force enumeration of the discretised posterior --- summing over all
$\ngrid^{\nsites}$ configurations on a short chain, sharing no code with the message passing ---
to $1.6\times10^{-14}$ on posterior means and $1.8\times10^{-15}$ on the log-evidence. Two
independent routes, not one self-consistent one.

\paragraph{What moment closure costs.} By Proposition~\ref{prop:closure-lmmse} the gap between
moment-projected and full sum-product is exactly the price of replacing the prior by its
covariance-matched Gaussian. Measured on the Laplace chain, the median relative score error falls
from $0.241$ at $t=0.05$ to $7.3\times10^{-6}$ at $t=3$ --- a factor of $3.3\times10^{4}$ across
the schedule, as noising Gaussianises the posterior and the innovation's shape stops mattering.

\section{Settled: learning the kernel}
\label{sec:rec-rung1}

\paragraph{Both optimisation routes reach the same place.} On the Gaussian kernel, gradient
ascent via Fisher's identity reaches the exact M-step's optimum to $0.021$ nats at its best step
size, with slightly better parameter errors ($0.0040$ against $0.0053$ on $\corr$). The
difference is robustness, not accuracy: EM's monotonicity violation is $0.0$ exactly and it needs
no step size, while gradient ascent at $\eta=2.0$ diverges ($\corr$ error $0.41$, monotonicity
violation $1641$) and at $\eta=8.0$ is unstable. That is the whole content of the comparison, and
it is why every rung above uses EM.

\paragraph{The gradient is exact.} Fisher's identity is verified against finite differences of
the discretised evidence to ${\sim}10^{-9}$. This check matters more than it looks: a fixed-point
residual near zero is evidence of a stationary point only if the derivative defining it is right,
and a sign error in exactly that derivative once survived every internal diagnostic
(\S\ref{sec:corr-gradient}).

\section{Settled: what the structure buys}
\label{sec:rec-efficiency}

Relative denoising error against grid BP under the true kernel, averaged over the noise schedule,
with \emph{both} arms' optimisation budgets chosen on a disjoint validation bundle
(\texttt{629985}):

\begin{center}\small
\begin{tabular}{@{}rccc@{}}
\toprule
$\nseq$ & network & EM--BP ($12$ free) & ratio \\
\midrule
32   & $0.3953 \pm 0.0023$ & $\mathbf{0.0546 \pm 0.0027}$ & $7.3$ \\
64   & $0.3285 \pm 0.0026$ & $\mathbf{0.0481 \pm 0.0041}$ & $7.9$ \\
128  & $0.2744 \pm 0.0012$ & $\mathbf{0.0328 \pm 0.0018}$ & $8.8$ \\
256  & $0.2268 \pm 0.0010$ & $\mathbf{0.0243 \pm 0.0014}$ & $10.1$ \\
512  & $0.1878 \pm 0.0006$ & $\mathbf{0.0214 \pm 0.0011}$ & $9.3$ \\
1024 & $0.1569 \pm 0.0006$ & $\mathbf{0.0162 \pm 0.0011}$ & $12.6$ \\
2048 & $0.1326 \pm 0.0005$ & $\mathbf{0.0134 \pm 0.0008}$ & $15.7$ \\
\bottomrule
\end{tabular}
\end{center}

All sixteen seeds, all seven doublings, $\pm$ one standard error \emph{aggregated per seed and then
across seeds} --- the twelve noise levels within a seed share one training set and one set of
fitted models, so they are not twelve independent observations. The network's
(parameterisation, training length) choice agrees with the test-set optimum in $82.2\%$ of the
$1{,}344$ cells and EM--BP's iteration count in $79.5\%$, so the selection is signal rather than
noise. EM--BP is more accurate in $1{,}342$ of $1{,}344$ cells; the two exceptions are both at
$\nseq=32$, $t=3.0$, at ratios $0.92$ and $0.96$. The network at $\nseq=2048$ ($0.133$) has not
reached the error the estimator attains at $\nseq=32$ ($0.055$), and no crossing is observed, so no
data-equivalence factor is quoted.

\begin{pitfall}
\textbf{The largest ratio in the table is the least trustworthy number in it.} At $\nseq=2048$ the
network selects the $20{,}000$-step cap in $33\%$ of cells and EM--BP selects the $400$-iteration
cap in $61\%$: both arms are asking for more budget than the grid offers, so $15.7$ is bounded by
the grid rather than by the estimators. Over the remaining six sizes the range is $7.3$--$12.6$.
Quote that range, not the endpoint.
\end{pitfall}

\paragraph{What the symmetric protocol cost, and why it was worth it.}
The previous version of this table (\texttt{629091}) selected EM--BP's iteration count on
validation while training the network to a fixed $20{,}000$ steps. Correcting that moved every
ratio down --- by $32\%$ at $\nseq=32$, $36\%$ at $128$, and only $5\%$ at $2048$ --- because early
stopping helps the network most where data is scarcest. Measured on each arm separately, running to
the configured budget instead of the selected one costs the network $+42\%$, $+52\%$, $+52\%$ at
$\nseq=32,64,128$ against EM--BP's $+39\%$, $+32\%$, $+14\%$. Pinning one arm's budget was
therefore not a neutral bookkeeping choice; it was worth up to a third of the headline number, in
our favour.

The ratio is not monotone and that shape should not be over-read. Both arms improve with $\nseq$;
the estimator starts so far ahead that the ratio is a quotient of two moving quantities, and
nothing here identifies a mechanism. A side observation, recorded but not claimed: the selected
network training length scales as $\nseq^{0.84}$ ($r=0.99$) against EM--BP's $\nseq^{0.53}$
($r=1.00$), so the network needs disproportionately more optimisation as data grows --- but both
slopes are lower bounds, since both arms hit their caps at $\nseq=2048$.

\begin{pitfall}
These sixteen seeds ran at $120$ EM iterations, because \codefile{exp\_07} hard-coded that budget
while \texttt{FROZEN.em\_max\_iters} said $400$ and nothing read it (\S\ref{sec:corr-capacity}).
The rerun at the configured budget is Bocconi array \texttt{628943}.

\textbf{The obvious prediction about that rerun was wrong, and the way it was wrong is the
interesting part.} The reasoning was: the shape settles at a median of $229$ updates, so a run
stopped at $120$ is short of convergence, so the estimator is handicapped and the ratios are a
lower bound. On the first six seeds the direction \emph{reverses with sample size}. Held-out score
error of the estimator, same six seeds, $120$ against $400$ iterations:

\begin{center}\small
\begin{tabular}{@{}rccc@{}}
\toprule
$\nseq$ & EM at $120$ & EM at $400$ & \\
\midrule
32   & $0.0585$ & $0.0673$ & worse by $15\%$ \\
64   & $0.0501$ & $0.0626$ & worse by $25\%$ \\
256  & $0.0283$ & $0.0295$ & worse by $4\%$ \\
512  & $0.0223$ & $0.0200$ & better by $10\%$ \\
2048 & $0.0158$ & $0.0125$ & better by $21\%$ \\
\bottomrule
\end{tabular}
\end{center}

\noindent
Running the optimiser \emph{longer} makes the estimator \emph{worse} at small $\nseq$ and better
at large $\nseq$, so no fixed budget is neutral between the two arms of the comparison.

\paragraph{The mechanism, measured rather than inferred.}
\label{sec:em-overfitting}
An instrumented run
(\codefile{experiments/exp\_29\_em\_overfitting.py}) records the training marginal log-evidence and
the held-out score error on the \emph{same} iterates, which is what separates overfitting from a
bad optimum: overfitting means the training objective keeps improving while generalisation turns
around, whereas a bad optimum would show both stalling.

\begin{center}\small
\begin{tabular}{@{}rccc@{}}
\toprule
$\nseq$ & training $\Lik$ monotone $\uparrow$ & held-out best at & cost of running to $400$ \\
\midrule
32   & yes & iteration $130$ & $+2.8\%$ \\
64   & yes & iteration $40$  & $+46.0\%$ \\
2048 & yes & iteration $250$ & $+8.6\%$ \\
\bottomrule
\end{tabular}
\end{center}

\noindent
The training log-evidence increases monotonically to the last iterate at every size --- as EM
guarantees it must --- while the held-out error has an \emph{interior minimum} at every size. That
is overfitting, and the diagnosis is now a measurement rather than a plausible story.

Note what the third row says: the interior optimum is at $250$ even at $\nseq=2048$. The effect is
\textbf{not} confined to small samples, and the earlier reading --- ``$120$ wins below the
crossover, $400$ wins above'' --- was an artefact of comparing two arbitrary points either side of
an optimum that moves. A cap of $400$ overshoots at every size tested.

\paragraph{Two consequences.} First, ``run it to convergence'' is not automatically the honest
choice. The iteration count is a hyperparameter with a bias--variance trade-off like any other, and
fixing it a priori quietly favours one arm; picking it on the test set would be exactly the sin
\S\ref{sec:corr-flat} records. It is now chosen on the validation bundle, per noise level,
symmetrically with the network's parameterisation, which is what \texttt{629091} runs.

Second --- and this is why it is in this chapter rather than a footnote --- this is the
memorisation--generalisation question of \S\ref{sec:memorisation} appearing in a model with twelve
free parameters and no network anywhere. The training objective and the generalisation objective
come apart, early stopping is doing implicit regularisation, and all of it is visible in a setting
where the target score is computable in closed form. That is precisely what this setting was built
for, and it is the most promising thing found in this batch.

\paragraph{Still not claimed.} That the optimal stopping point scales with $\nseq$ in any
particular way. The three sizes here give $130$, $40$ and $250$, which is not monotone, and three
points chosen for a diagnostic are not a curve. Reading a scaling law off them would repeat
\S\ref{sec:corr-flat} exactly.
\end{pitfall}

\section{Settled: structure the marginals cannot see}

Chapter~\ref{ch:ring} in full. The result in one line: across $1{,}344$ cells the per-frame arm's
estimate of $p(\psi)$ moves $0.0000$ in total variation from its uniform initialisation, and
across the $480$ cells of the single-rotation variant its profile likelihood in $\psi$ has range
$0.00\times10^{0}$ nats --- a machine-precision zero, which is what
Theorem~\ref{thm:ring-blind} requires. The joint arm recovers the rotation to $0.15^\circ$ at
$t=0.05$, degrading to $7.09^\circ$ at $t=1.43$ and returning to the null by $t=3$.

\section{Settled: where the Markov assumption fails}
\label{sec:rec-nonmarkov}

Two contamination mechanisms, both with an exact reference, from the complete ten-task run
\texttt{627165}. \texttt{ratio\_to\_em} is the baseline's relative score error over EM--BP's, so
$>1$ means the estimator wins.

\begin{center}\small
\begin{tabular}{@{}lcccccc@{}}
\toprule
& \multicolumn{5}{c}{global latent, rank-one, strength $\beta$} \\
\cmidrule(lr){2-6}
& $0.00$ & $0.10$ & $0.25$ & $0.50$ & $1.00$ \\
\midrule
Gaussian, vs CNN & $18.3$ & $17.2$ & $6.99$ & $3.36$ & $2.22$ \\
Laplace, vs CNN  & $9.90$ & $10.6$ & $6.50$ & $3.05$ & $2.12$ \\
\midrule
& \multicolumn{5}{c}{long-range precision coupling, strength $\gamma$} \\
\cmidrule(lr){2-6}
& $0.00$ & $0.05$ & $0.10$ & $0.20$ & $0.40$ \\
\midrule
Gaussian, vs CNN & $24.1$ & $1.19$ & $\mathbf{0.98}$ & $0.86$ & $0.77$ \\
Gaussian, vs MLP & $37.0$ & $1.79$ & $1.23$ & $0.96$ & $0.83$ \\
\bottomrule
\end{tabular}
\end{center}

The estimator survives rank-one contamination essentially intact --- still winning at $\beta=1$,
where half the marginal variance is a shared constant --- and \textbf{inverts at
$\gamma\approx0.10$} under long-range coupling. The mechanism is exactly what the structure
predicts: a global latent makes the score residual rank one, which a chain absorbs; long-range
coupling is not low-rank and cannot be represented. This is the honest scope statement, and it is
backed by a complete run.

\begin{pitfall}
These runs predate \codefile{frozen\_config.py} and are \emph{not} at the frozen settings. They
are reported here, in the development document, and are deliberately not quoted in the note or
the workshop paper. The qualitative conclusion --- rank-one survivable, long-range not --- is
robust to the settings; the specific ratios are not to be pooled with
\S\ref{sec:rec-efficiency}.
\end{pitfall}

\section{In progress}
\label{sec:pending}

Four of the five reruns listed in the previous version have landed and moved into the sections
above; what follows is what they said and what is still out.

\paragraph{Landed: the parametric rate} (\texttt{628576}). At four replicates the fitted
$\log$--$\log$ slopes were $-0.196$ for $\corr$ and $-0.700$ for $\ivar$ against a predicted
$-0.500$, with a non-monotone $\corr$ curve --- not a confirmation of $\nseq^{-1/2}$, and the
claim was pulled from the note rather than caveated. At sixteen the variance slope is $-0.519$
($r=-0.99$, monotone), which is the rate; $\corr$ is monotone but shallower at $-0.273$. Half the
claim is restored, half is not, and the note says so rather than averaging them into a verdict.

\paragraph{Landed: clean against channel} (\texttt{628600}). Two slownesses, and they are
independent. \emph{Optimisation}: coordinates of one kernel settle at very different rates ---
median $80$ updates for $\corr$, $68$ for $\ivar$, $229$ for the excess kurtosis
(\texttt{628601}, below). \emph{Statistics}: fitting through the channel costs rate, not merely a
constant.

\begin{center}\small
\begin{tabular}{@{}lcc@{}}
\toprule
$\log$--$\log$ slope, sixteen replicates & $\corr$ & $\ivar$ \\
\midrule
clean transition pairs & $-0.52$ & $-0.50$ \\
through the channel    & $-0.44$ & $-0.36$ \\
\bottomrule
\end{tabular}
\end{center}

\noindent
The clean arm sits on the parametric rate on both coordinates, which is the check that the
estimator is not itself the bottleneck. It is also flatly incompatible with the withdrawn
discretisation-floor reading of \S\ref{sec:corr-flat}: the clean arm falls by a factor of $3.4$ on
$\corr$ and $4.3$ on $\ivar$ across the range, so there is no floor to see.

\paragraph{Landed: shape convergence at sixteen seeds} (\texttt{628601}). The settling table above.
Zero monotonicity violations across all sixteen traces. The kurtosis figure has a long tail --- up
to $638$ updates in the slowest seed --- which is what makes a $400$ cap a decision rather than a
formality.

\paragraph{Landed: sample efficiency, both budgets selected} (\texttt{629985}, sixteen seeds).
\S\ref{sec:rec-efficiency}. Supersedes \texttt{628571} (fixed budgets both sides) and
\texttt{629091} (EM-BP's budget selected, the network's pinned). This is the version the paper
carries, with the $\nseq=2048$ caveat recorded there.

\paragraph{Landed, partly: the confining potential} (Rung 4a, \texttt{629193}).
\label{sec:rung4a-status}
Three attempts, and the third has now landed in full, so this is worth setting out in order.

\paragraph{Result, complete (\texttt{629962}: $480$ cells, eight seeds $\times$ five sizes
$\times$ six noise levels, balanced).} The joint arm recovers the well width better than the
marginal arm --- median relative error $0.277$ against $0.404$ --- and its accuracy improves
monotonically with the number of trajectories:

\begin{center}\small
\begin{tabular}{@{}lccccc@{}}
\toprule
$\nseq$ & $128$ & $512$ & $1024$ & $2048$ & $4096$ \\
\midrule
median $|\hat\lambda-\lambda|/\lambda$ & $0.656$ & $0.367$ & $0.260$ & $0.192$ & $0.161$ \\
\bottomrule
\end{tabular}
\end{center}

\noindent
The convergence fix is what made this readable. Under the previous absolute tolerance the
marginal arm hit its iteration cap in $87.8\%$ of cells against the joint arm's $56.7\%$, so the
two arms were being compared at different amounts of optimisation; only $13$ of $400$ pairs had
both converged. With the per-sequence tolerance and a raised cap the rates are $44.2\%$ and
$21.2\%$, and the comparison is between estimators rather than between budgets.

\paragraph{Where the advantage lives, and where it reverses.} Resolved by noise level, the
ordering is not schedule-wide and the crossover is sharp:

\begin{center}\small
\begin{tabular}{@{}lcccccc@{}}
\toprule
$t$ & $0.050$ & $0.105$ & $0.222$ & $0.467$ & $0.982$ & $2.068$ \\
\midrule
joint    & $\mathbf{0.086}$ & $\mathbf{0.137}$ & $\mathbf{0.127}$ & $\mathbf{0.359}$ & $0.860$ & $0.936$ \\
marginal & $0.141$ & $0.170$ & $0.447$ & $0.410$ & $\mathbf{0.754}$ & $\mathbf{0.838}$ \\
\bottomrule
\end{tabular}
\end{center}

\begin{pitfall}
Two caveats travel with the headline, and the schedule-wide median hides both.

The joint arm wins at $t \leq 0.467$ and \emph{loses} at $t \geq 0.982$. Past the crossover
the channel has destroyed most of the rotation information, both arms are near the null, and
the joint likelihood's extra freedom costs more than it buys. Quoting $0.277$ against $0.404$
without the breakdown would imply a uniform advantage that the data does not show.

And the errors are large in absolute terms. At the largest sample size the joint arm is still
$16\%$ off the true well width. What this rung supports is the \emph{ordering} at low noise and
the \emph{scaling} in $\nseq$ --- not a precision claim about $\lambda$.
\end{pitfall}

\texttt{628434\_5} was killed by the six-hour wall having written nothing --- $960$ gradient fits at
${\sim}0.2$\,s per step is roughly ten hours, and the experiment writes once at the end. Sharding by
seed fixed that. \texttt{628942} then completed and produced nothing usable, for a reason that had
nothing to do with the queue: \texttt{part\_potential} generated two species at $\pm\omega$ while
\texttt{fit\_potential} scores every trajectory at the single scalar $\psi$ it is handed, so half the
data was evaluated under the wrong rotation. The estimator's escape was to collapse the ring. At
$r_*\to0$ the ring is rotation-invariant and so pays nothing for a wrong $\psi$; $r_*$ pinned to the
lower clip in $51.2\%$ of cells, and the sweep reported the \emph{marginal} arm beating the joint
arm, which is the reverse of what this rung exists to show.

\texttt{629193} (one species, $600$ steps, sixteen seeds, $960$ cells) repairs that:

\begin{center}\small
\begin{tabular}{lccc}
\toprule
 & $\lambda$ rel.\ err.\ (med.) & $r_*$ abs.\ err.\ (med.) & pinned at clip \\
\midrule
joint, two species (\texttt{628942}) & $1.844$ & $0.950$ & $51.2\%$ \\
joint, one species (\texttt{629193}) & $0.229$ & $0.020$ & $4.8\%$ \\
marginal, one species                & $0.429$ & $0.039$ & $4.4\%$ \\
\bottomrule
\end{tabular}
\end{center}

\noindent The joint arm's error falls monotonically with sample size --- $0.618$, $0.336$, $0.222$,
$0.181$, $0.118$ at $n = 128$ to $4096$ --- and that trend is reported here because it survives the
caveat below.

\begin{pitfall}
The arm \emph{comparison} from \texttt{629193} is not reported, because the two arms were not held to
the same convergence standard. At $t\le1$ the marginal arm hit the $600$-step cap in $95\%$ of cells
against the joint arm's $57\%$, so ``joint beats marginal'' partly means ``joint converged and
marginal did not''. The obvious control does not rescue it: only $13$ of $400$ pairs have both arms
converged.

The cause was the stopping rule. \texttt{plateau\_tol} was an \emph{absolute} tolerance on an
objective that sums over sequences, hence eight times stricter at $n{=}4096$ than at $n{=}512$: a
convergence criterion varying along the very axis this rung reports its scaling on. It is now
per-sequence, pinned by
\texttt{tests/test\_ring.py::test\_stopping\_rule\_does\_not\_depend\_on\_sample\_size}, which
replicates one dataset fourfold and asserts the fit stops at the same step.
\end{pitfall}

\paragraph{But more optimisation is not more accuracy, and that is the deeper problem.}
\label{sec:rung4a-nonmonotone}
Twelve cells were refit to $2400$ steps with the plateau test disabled, to see what the budget was
costing. The first cell inspected showed exactly the expected picture --- at $n{=}512$,
$t{=}0.2216$, the joint arm is identical to eight significant figures from step $600$ onward while
the marginal arm's $\lambda$ error falls from $0.482$ to $0.372$ to $0.335$ --- and it would have
been easy to stop there and report that the joint arm was converged and the marginal arm cut short.
Across all twelve cells that is not what happens:

\begin{center}\small
\begin{tabular}{lccc c}
\toprule
cell & \multicolumn{3}{c}{$\lambda$ rel.\ err.\ at $600 / 1200 / 2400$} & log-lik.\ monotone \\
\midrule
$n{=}512$, $t{=}0.2216$, joint    & $0.078$ & $0.078$ & $0.078$ & yes \\
$n{=}512$, $t{=}0.4665$, joint    & $0.029$ & $0.057$ & $0.060$ & yes \\
$n{=}512$, $t{=}0.2216$, marginal & $0.177$ & $0.066$ & $0.207$ & yes \\
$n{=}2048$, $t{=}0.2216$, marginal& $0.265$ & $0.075$ & $0.149$ & yes \\
\bottomrule
\end{tabular}
\end{center}

\noindent The objective rises monotonically in all twelve cells --- backtracking guarantees it ---
while the parameter error is non-monotone in four, and in two joint-arm cells the estimate at $2400$
steps is \emph{worse} than at $600$. Ascending the likelihood harder does not monotonically improve
recovery, because the finite-sample maximiser is not at the truth. This is the same phenomenon as
the EM overfitting of \S\ref{sec:em-overfitting}, in a different estimator: the optimisation budget
is a regularisation knob, not a convergence detail.

So \texttt{629962} should be read for what it is. Holding both arms to one $n$-independent
criterion removes a defect that was certainly there; it does not deliver two converged estimators,
and the honest fix --- selecting each arm's stopping point on held-out data, as
\S\ref{sec:rec-efficiency} now does for EM-BP and the network --- is not implemented for this rung.
Any joint-versus-marginal margin from \texttt{629962} carries that caveat.

\paragraph{Still out.}
\begin{itemize}\itemsep3pt
\item \textbf{Rung 4a at a matched convergence budget} (\texttt{629962}), eight seeds split by
      sample size into sixteen shards, at a $3000$-step cap under the per-sequence tolerance. Until
      it lands, the joint-versus-marginal margin above stays out of every document.
\end{itemize}

\section{Withdrawn, and not replaced}

Recorded in full in Chapter~\ref{ch:corrections}; listed here so that nothing in this chapter is
read as still standing.

\begin{itemize}\itemsep3pt
\item \textbf{The capacity attribution.} That enlarging the innovation mixture buys generative
      fidelity. Roughly $96\%$ of the apparent effect was convergence rate at a fixed iteration
      budget. The rerun that would settle it lost all but one of its array tasks, so one cell
      survives and the claim stands withdrawn and unreplaced (\S\ref{sec:corr-capacity}).
\item \textbf{The discretisation floor.} That the clean-arm variance error was floored by the
      grid. It was under-replicated, not floored, and the grid contributes a constant
      $4.4\times10^{-4}$ --- an order of magnitude too small to produce the effect
      (\S\ref{sec:corr-flat}).
\item \textbf{Everything downstream of the superseded generation protocol}, including the
      four-family comparison at matched covariance (\S\ref{sec:corr-generation}).
\end{itemize}

\section{Not started}

The architectural question. Sum-product on a chain suggests a network that is narrow but as deep
as twice the chain, unlike the tree-structured case where the natural depth is the tree's
\citep{mei2024unet,garnierbrun2024transformers}. Which architectures exploit Markovianity, and at
what depth, is what this setting was built to measure, and is the obvious continuation. The
receptive-field sweep that would begin it exists (\codefile{exp\_12}) but selected its radius on
the evaluation set, so its numbers are exploratory and are not reported anywhere.


\section{Corrections: what was wrong, and how it was caught}
\label{part:audit-corrections}

Each entry below is a claim that was made and then withdrawn or narrowed. They
are kept because the mechanism that produced each one is still live in the
code, and because the pattern across them is itself a finding.

This chapter exists so that the note and the workshop paper do not have to carry it. A reader of
those documents should meet claims, not retractions of claims; a reader of this one should be
able to see what was believed, why it was wrong, and --- the part that matters for the next
mistake --- what diagnostic caught it.

Every entry follows the same shape: the claim as it was made, the correction, and the check that
now exists so it cannot recur silently. Several of these were live in circulated drafts.

\subsection{The iteration budget that manufactured a capacity effect}
\label{sec:corr-capacity}

\paragraph{The claim.} At $\ncomp=4$ the estimator was an order of magnitude better pointwise
than the convolutional baseline and yet generated a worse residual law; varying $\ncomp$ improved
both axes monotonically, which identified the deficit as a \textbf{capacity} limit of the
innovation mixture.

\paragraph{What was wrong.} Every run behind that account used \texttt{em\_iters}${}=40$. EM on
this kernel converges on $\corr$ far faster than on the innovation \emph{shape}: at $\corr=0.85$,
$\nseq=512$, $\ncomp=8$ the fitted excess kurtosis reads $0.84$ at 30 iterations, $1.85$ at 60,
$2.30$ at 120 and $2.29$ at 400, while $\corr$ is flat from iteration ${\sim}25$. Monitoring
$\corr$ --- which is the natural thing to plot --- gives a converged-looking trace over a kernel
that is nowhere near converged. Enlarging the mixture at a fixed iteration budget buys
convergence \emph{rate}, not representational capacity, and roughly $96\%$ of the apparent
capacity benefit was that.

\paragraph{What replaced it.} Nothing yet, and this is worth stating plainly. The rerun that
would settle it was submitted as Bocconi array \texttt{627164}, and \textbf{only array task 2
ever ran} --- \texttt{sacct} lists \texttt{627164\_2} and nothing else. One cell survives,
$\ncomp=2$ at one seed, in which the estimator is consistent with the reference
($1.837\pm0.061$ against $1.894\pm0.096$, a gap of $1.20\sigma$). That is suggestive that the
capacity reading was indeed a convergence artefact --- a two-component mixture would not have
sufficed under the old account --- but one cell at one capacity cannot replace a withdrawn
claim. \textbf{The capacity attribution remains withdrawn and unreplaced.}

\paragraph{The check that was supposed to exist, and did not.} This section previously claimed
that \codefile{experiments/frozen\_config.py} defines a stopping rule on the \emph{shape} --- EM
stopping when $|\Delta\kappa_4| < 10^{-3}$, with a cap of 400 iterations. \textbf{That was false
when written and stayed false for a month.} The config declares \texttt{em\_shape\_tol} and
\texttt{em\_max\_iters}, but a search of the tree finds \emph{no reader of either}: the rule
actually implemented in \coderef{src/em.py}{fit\_em} is a relative tolerance of $10^{-9}$ on the
log-evidence, and the cap is whatever the caller passes. \codefile{exp\_07} --- the headline
experiment --- passed a literal $120$, in four places.

This is the replicate defect of \S\ref{sec:corr-replicates} in another costume, and it is worse
in one respect: there the compendium recorded the gap once it was found, whereas here the
compendium asserted a safeguard that did not exist. A document describing its own checks is
evidence about those checks only if the description is itself checked.

Two mitigating facts, neither of which excuses it. The tolerance that \emph{is} implemented is far
stricter than the one advertised: replayed against the 800-update traces of
\codefile{exp\_27}, it fires in one of sixteen seeds, so in practice runs go to the cap rather
than stopping early on a false plateau. And the direction of the exp\_07 error is conservative ---
an under-iterated EM makes the estimator look worse, so the efficiency ratio was understated, not
inflated.

\paragraph{The check that now exists, and what it immediately found.} \codefile{exp\_07} reads
\texttt{FROZEN.em\_max\_iters}, so the budget is the configured one.
\codefile{tools/check\_replicates.py} is extended from replicate counts to iteration budgets, and
--- this is the part that matters --- from settings dictionaries to \emph{call sites}, since
\texttt{n\_iters=120} is a keyword argument and the dict-only version could never have seen it.

Turned on, it reports sixteen fixed budgets across three paper experiments. Sorting them by
whether the run reports a coordinate the budget can actually reach:

\begin{center}\small
\begin{tabular}{@{}llcl@{}}
\toprule
experiment & sites & budget & reports \\
\midrule
\codefile{exp\_06} & 6 & 80--120 & $\corr$, $\ivar$ only --- settle at $80$, $68$: adequate \\
\codefile{exp\_06} & 6 & 120 & innovation \emph{moments} --- kurtosis settles at $229$ \\
\codefile{exp\_18} & 1 & 60 & innovation moments \\
\codefile{exp\_18} & 1 & 60 & $\corr$ only: adequate \\
\bottomrule
\end{tabular}
\end{center}

So the defect is not confined to \codefile{exp\_07}. Eight of the sixteen budgets are defensible
because the coordinates they report settle well inside them; \textbf{eight are not}, and every
number downstream of those is a kernel stopped roughly half-way to its shape. None of them is
quoted in the note or the workshop paper --- the shape-dependent claims there come from
\codefile{exp\_07}, which is being rerun --- but the fitted-density material in this document
does depend on them, and is to be read as provisional until they are rerun at the frozen budget.

\begin{pitfall}
An iteration budget is not a global property of a run. It is adequate or not \emph{relative to the
slowest coordinate that run reports}, and the same 120 iterations are ample for $\corr$ and
insufficient for a fourth cumulant. This is why \codefile{exp\_27} exists and why its settling
table is a prerequisite for reading any other experiment here, rather than a curiosity about
convergence.
\end{pitfall}

\codefile{experiments/exp\_27\_shape\_convergence.py} records the whole diagnostic set at every
evaluated state, which is what the original experiment never did.

\subsection{The replicate counts the frozen configuration could not reach}
\label{sec:corr-replicates}

\paragraph{The claim.} That a frozen configuration file, imported by every experiment, was
sufficient to guarantee that every number in the paper came from one protocol.

\paragraph{What was wrong.} \codefile{frozen\_config.py} exposes \texttt{n\_seeds}. Four
experiments carried their own replicate knobs and read none of it:

\begin{center}\small
\begin{tabular}{@{}llrr@{}}
\toprule
experiment & knob & was & protocol \\
\midrule
\codefile{exp\_06} & \texttt{n\_rep} & 10 & 16 \\
\codefile{exp\_06} & \texttt{n\_rep\_rate} & 4 & 16 \\
\codefile{exp\_06} & \texttt{n\_rep\_clean} & 3 & 16 \\
\codefile{exp\_27} & \texttt{n\_seeds} & 8 & 16 \\
\codefile{exp\_07} & \texttt{seed} & one per run, four run & 16 \\
\bottomrule
\end{tabular}
\end{center}

\noindent
All five runs reported \texttt{is\_frozen: true}, because the provenance check compares the
config against itself and cannot see a knob the experiment never asked about. \emph{A
configuration that cannot reach the parameter that matters is not freezing anything.}

The consequences were real and different in each case. The rate experiment at four replicates
gives fitted $\log$--$\log$ slopes of $-0.196$ for $\corr$ and $-0.700$ for $\ivar$ against a
predicted $-0.500$, with a non-monotone $\corr$ curve --- that is not a confirmation of the
parametric rate, and the claim was pulled from the note rather than reported with a caveat. The
clean-versus-noised comparison at three replicates is the same count that produced the withdrawn
flatness reading of \S\ref{sec:corr-flat}.

\paragraph{The check that now exists.} \codefile{tools/check\_replicates.py} parses each
experiment's settings dictionaries and fails on any literal replicate count in a paper
experiment, allowing deviations only with an explicit \texttt{\# frozen-exempt: <reason>}
comment. It is worth recording that this check was \emph{first written as a shell grep}, and the
grep missed \codefile{exp\_27} entirely --- a line-based search cannot tell a \texttt{quick} smoke
dictionary from a \texttt{full} one. The AST version found it immediately. A check that is
approximately right about where to look is not a check.

\subsection{The flat curve that was under-replicated, not floored}
\label{sec:corr-flat}

\paragraph{The claim.} On clean chains the innovation-variance error was flat in $\nseq$
($0.0041, 0.0039, 0.0059, 0.0046, 0.0040$ from $\nseq=64$ to $1024$), and the flatness was a
discretisation floor imposed by the grid.

\paragraph{What was wrong.} Both halves. At three replicates an RMSE is estimated from three
numbers and its own sampling error is comparable to the effect being read off it; at sixteen the
curve is not flat, running $0.0136, 0.0073, 0.0056, 0.0042, 0.0032$, close to the $\nseq^{-1/2}$
reference. The $\nseq=64$ entry alone moves by a factor of three between the two replication
levels.

Nor was the grid ever a plausible culprit. Removing it entirely --- exact clean-data OLS on the
raw transition pairs, no binning --- gives $0.0135, 0.0071, 0.0057, 0.0043, 0.0032$, statistically
indistinguishable from the binned fit, with a paired grid-minus-raw discrepancy of a constant
$4.4\times10^{-4}$ across all budgets. A quantity that small cannot produce a floor at
$4\times10^{-3}$.

\paragraph{The lesson, which generalises.} The failure mode was reading a \emph{mechanism} into a
pattern that was noise, and then reaching for the most technically interesting available
explanation. The discretisation floor was a satisfying story; it was also never tested against
the alternative that three points are not enough to see a trend.

\subsection{Two wrong ways to compare a generated law}
\label{sec:corr-generation}

\paragraph{The claim.} Successive versions compared the generated sample against clean data, and
then applied a denoising read-out to it.

\paragraph{What was wrong.} Both, in opposite directions. Reverse integration stops at
$\tmin>0$, so what each estimator produces is a sample of $P_{\tmin}$, not of $\Pz$; comparing it
against clean data measures the stopping point rather than the score. Concretely, independent
Gaussian noise of variance $v$ on an innovation of variance $\ivar$ and excess kurtosis $\kappa$
gives excess kurtosis $\kappa(\ivar/(\ivar+v))^2$, since fourth cumulants add and the variance
grows. The second version applied a denoising read-out instead, which inverted the bias through
shrinkage.

\paragraph{What replaced it.} Draw the reference from $\Pz$ and noise it \emph{forward} to the
same $\tmin$, so both sides are samples of the same law and the target follows in closed form. At
$\corr=0.85$, $\tmin=0.02$ the residual innovation noise is $v=\Dt(1+\corr^2)$ and the predicted
target is $1.9098$ against a clean $3.0$. The closed-form target is what identified both errors:
neither was visible as an anomaly, and both produced plausible numbers.

\subsection{A sign error in a gradient that a fixed point could not see}
\label{sec:corr-gradient}

\paragraph{What was wrong.} \coderef{src/kernels.py}{MixtureInnovationKernel.grad\_log\_transition\_matrix}
carried a sign error in the $\corr$ derivative.

\paragraph{Why it mattered more than it looks.} A fixed-point residual near zero is evidence of a
stationary point only if the derivative defining it is right. Every convergence diagnostic
downstream of that function inherited the error, so the runs looked converged for a reason
unrelated to whether they were.

\paragraph{The check that now exists.} Fisher's identity is verified against finite differences
of the discretised evidence to ${\sim}10^{-9}$ --- a check that compares two independent routes to
the same quantity, rather than one that asks a quantity whether it is happy with itself.

\subsection{The ring normaliser, twice}
\label{sec:corr-ring}

Recorded in full in \S\ref{sec:blindness} and its surrounding boxes; summarised here because it
is the most recent and the most instructive.

The ring density is written unnormalised, and its normaliser depends on the parameters of the
confining well. Omitting it is harmless while only $\psi$ is being estimated --- the term is
constant and cancels in every responsibility --- and silently fatal the moment $\lambda$ is
estimated, because a wider well contains more mass and the likelihood grows without bound.

Two things are worth extracting. First, the bug was \emph{introduced by a correct-looking
port}: the ported code agreed with its audited original bit for bit, because the original only
ever compared $\psi$ at fixed $\lambda$ and never needed the term. Agreement with a reference
only certifies the cases the reference was used for. Second, fixing it in
\coderef{src/ring.py}{log\_pt\_conditional} left it unfixed in
\coderef{src/ring.py}{log\_pt\_marginal}, where the count differs --- once per frame rather than
once --- and that second instance was caught not by a test but by an experiment diverging to
$\hat\lambda = 3.5\times10^{6}$. The test suite now pins both functions and both directions.

\subsection{What these have in common}

Five of the six were failures of \emph{measurement}, not of derivation. No theorem was wrong. In
each case a quantity was computed correctly and then read as evidence for something it could not
support: an iteration budget read as a capacity limit, three replicates read as a floor, a
stopping point read as a score error, a fixed point read as convergence, a constant read as
absent.

Two habits follow, and they are the reason this project's checks look the way they do.

\paragraph{Prefer two independent routes to one self-consistent one.} The gradient sign error
survived every internal diagnostic and died immediately against finite differences. The
generation comparison survived two revisions and died against a closed-form target. Grid BP is
checked against a closed form \emph{and} against brute-force enumeration sharing no code with it.

\paragraph{Make the failing direction fail.} It is not enough for a test to pass when the code is
right; where a bug's signature is known, assert it. \codefile{tests/test\_ring.py} asserts that
the profile likelihood peaks at the truth \emph{and} that it becomes monotone when the normaliser
is removed, so deleting the term cannot make the suite greener. Most of the entries above would
have been caught a good deal sooner by a test of that shape.
